\scanpage{90}

que \deleted{\unsure{pour}} $(j,k\in J\text{ et }j\leq i\leq k)\Longrightarrow i\in J$. On dit aussi alors, que $J$ est une \uline{bande} de $I$. C'est \deleted{\unsure{une propriété}} une propriété sym. par passage de $I$ à $I^\circ$.
\inserted{Toute intersection de bandes est une bande. La plus petite bande contenant $J$ est \deleted{\unsure{notée}}
\[
(84)\qquad\underbrace{\bar J\cap\widetilde J}_{\substack{\Vert\,\mathrm{déf}\\B(J)}}=\{i\in I\mid\exists j,k\in J\quad j\leq i\leq k\}
\]}

Soient $J,K$ deux \deleted{\unsure{bandes}} \inserted{parties} de $I$, on écrit
\[
(85)\qquad J\prec K\quad\overset{\mathrm{déf}}{\Longleftrightarrow}\quad J\subset\bar K,\ K\subset\widetilde J
\]
(relation de \uline{domination}). C'est une relation de préordre sur l'ens $\mathcal P(I)$ des parties de $I$, la relation d'équivalence associée est
\[
(86)\qquad\begin{aligned}
J\sim K&\Longleftrightarrow J\subset\underbrace{\bar K\cap\widetilde K}_{B(K)},\ K\subset\underbrace{\bar J\cap\widetilde J}_{B(J)}\\
&\Longleftrightarrow B(J)=B(K)
\end{aligned}
\]
L'ens ordonné quotient est l'ens des \uline{bandes} de $I$, ordonné par la relation de \deleted{\unsure{l'inclusion}} domination.

Supposons maintenant que $J$ et $K$ soient des bandes, \inserted{avec $J\prec K$,} soit $I'$ la bande engendrée par $J\cup K$, on a
\[
(87)\qquad\begin{aligned}
I'&=\{i\in I\mid\exists j\in J,k\in K\text{ t.q. }j\leq i\leq k\}\\
&=\widetilde J\cap\bar K
\end{aligned}
\]
et $J$ est fermé dans $I'$, $K$ est ouvert dans $I'$ enfin \deleted{\unsure{et}} $K$ est dense dans $I'$ et $J$ codense dans $I'$ i.e. $I'\subset\bar K$ et $I'\subset\widetilde J$ \deleted{\unsure{et}} i.e. $I'\subset\bar K\cap\widetilde J$ (on a $=$ \uline{égalité}!). Inversement si $\exists$ bande $I'$ tel que $J$ soit un fermé codense dans $I'$, $K$ un ouvert dense dans $I'$, alors $J$ et $K$ sont des bandes, $J\prec K$, et \unsure{ainsi} $I'=\widetilde J\cap\bar K$.
\margin{89}
